% Gauss, Werke X/1 (1917), native leaves 35–39, printed 28–32.
\SourceLeaf{035}\NoSourceHead
\SourceGroup{III}
\HistoricTitle{FÜNFTER BEWEIS DES FUNDAMENTALSATZES BEI DEN QUADRATISCHEN RESTEN.}
\HistoricCitation{[Aus Handbuch 21, Bg. Aufsätze, Notizen und Rechnungen zur Mathematik und Astronomie gehörig. Angefangen September 1813, S. 4—5.]}
\par\indent I) Es sei $p$ eine (ungerade) Primzahl, $a$ eine Radix primitiva für den Modulus $p$, $x$ eine unbestimmt bleibende Grösse. Wir bezeichnen durch $fx$ die Function
\[
[1+]x+x^a+x^{aa}+x^{a^3}+x^{a^4}+x^{a^5}+\text{etc.}+x^{a^{p-2}}.
\]
\noindent Man sieht leicht, dass
\[
fx-([1+]x+xx+x^3+\text{etc.}+x^{p-1})
\]
\noindent durch $1-x^p$ theilbar seyn müsse, da die Zahlen $1,a,aa,a^3\ldots,a^{p-2}$ den Zahlen $1,\allowbreak2,\allowbreak3,\allowbreak4,\allowbreak\ldots,\allowbreak p-1$ nach dem Modulus $p$ congruent sind (ohne Rücksicht auf die Ordnung). Es ist also $fx$ durch
\[
1+x+xx+\text{etc.}+x^{p-1}=\tfrac{1-x^p}{1-x}
\]
\noindent theilbar. Es wird also auch allgemein, wenn $n$ irgend eine ganze Zahl bedeutet, $f(x^n)$ durch $\frac{1-x^{np}}{1-x^n}$ theilbar seyn. Nun ist aber leicht zu beweisen, dass
\[
\tfrac{1-x^{np}}{1-x^n}\text{ durch }\tfrac{1-x^p}{1-x}
\]
\noindent theilbar sei, wenn $n$ durch $p$ nicht theilbar ist, in diesem Fall wird also $f(x^n)$ auch durch $\frac{1-x^p}{1-x}$ theilbar seyn. Hingegen sieht man leicht, dass wenn $n$
\clearpage
\SourceLeaf{036}\ArithmeticHead{29}{R}{FÜNFTER BEWEIS DES FUNDAMENTALSATZES.}
\noindent durch $p$ theilbar ist, $f(x^n)-p$ durch $1-x^p$ also auch durch $\frac{1-x^p}{1-x}$ theilbar seyn w[ird].

II) Man bezeichne ferner
\[
x-x^a+x^{aa}-x^{a^3}+\text{etc.}-x^{a^{p-2}}
\]
\noindent durch $\xi$. Alsdann ist
\[
+x\xi-(x^2-x^{a+1}+x^{aa+1}-x^{a^3+1}+\text{etc.}-x^{a^{p-2}+1})=0
\]
\noindent und die Grössen
\begingroup\small
\[
\begin{aligned}
-x^a\xi&-(x^{2a}-x^{aa+a}+x^{a^3+a}-x^{a^4+a}+\text{etc.}-x^{a^{p-1}+a}),\\
+x^{aa}\xi&-(x^{2aa}-x^{a^3+aa}+x^{a^4+aa}-x^{a^5+aa}+\text{etc.}-x^{a^p+aa}),\\
-x^{a^3}\xi&-(x^{2a^3}-x^{a^4+a^3}+x^{a^5+a^3}-x^{a^6+a^3}+\text{etc.}-x^{a^{p+1}+a^3}),\\
+x^{a^4}\xi&-(x^{2a^4}-x^{a^5+a^4}+x^{a^6+a^4}-x^{a^7+a^4}+\text{etc.}-x^{a^{p+2}+a^4})
\end{aligned}
\]
\begin{center}u. s. w. bis zu\end{center}
\[
-x^{a^{p-2}}\xi-(x^{2a^{p-2}}-x^{a^{p-1}+a^{p-2}}+x^{a^p+a^{p-2}}-x^{a^{p+1}+a^{p-2}}+\text{etc.}-x^{a^{2p-4}+a^{p-2}})
\]
\endgroup
\noindent jede einzeln durch $1-x^p$ theilbar. Addirt man alles, so ist folglich klar, dass
\[
\xi\xi-f(x^2)+f(x^{a+1})-f(x^{aa+1})+f(x^{a^3+1})-\text{etc.}+f(x^{a^{p-2}+1})
\]
\noindent durch $1-x^p$ theilbar seyn werde, also auch durch $\frac{1-x^p}{1-x}$. Nun ist aber bekanntlich unter den Zahlen
\[
2,\ a+1,\ aa+1,\ a^3+1\text{ etc. bis }a^{p-2}+1
\]
\noindent nur eine einzige nemlich $a^{\frac12(p-1)}+1$ durch $p$ theilbar: dem vorhergehenden Theorem zufolge wird also
\[
\xi\xi\mp f(x^{a^{\frac12(p-1)}+1})
\]
\noindent und folglich $\xi\xi\mp p$ durch $\frac{1-x^p}{1-x}$ theilbar seyn, wo das obere oder untere Zeichen gilt, je nachdem $\tfrac12(p-1)$ gerade oder ungerade, d. i. je nachdem $p$ von der Form $4m+1$ oder $4m-1$ ist.
\clearpage

\SourceLeaf{037}\ArithmeticHead{30}{L}{ARITHMETIK.\quad NACHLASS.}
\par\indent III) Es sei nun $q$ irgend eine ungerade (positive) Zahl also $\tfrac12(q-1)$ eine ganze Zahl. Es wird folglich
\[
(\xi\xi)^{\frac12(q-1)}-(\pm p)^{\frac12(q-1)}
\]
\noindent durch $\xi\xi-(\pm p)$, also auch durch $\frac{1-x^p}{1-x}$ theilbar seyn; jene erstere Grösse kann man auch so schreiben:
\[
\xi^{q-1}\mp p^{\frac12(q-1)},
\]
\noindent wo das obere Zeichen gelten wird, wenn von dem Zahlen $p,q$ wenigstens eine von der Form $4m+1$ ist, das untere hingegen, wenn beide von der Form $4m+3$ sind.

IV) Nehmen wir jetzt noch an, dass auch $q$ eine von $p$ verschiedene Primzahl ist, so folgt aus dem Theorem, welches in den \textit{Disqu. Ar.} p. 46\begingroup\renewcommand{\thefootnote}{[*)}\footnote{Werke I, S. 42.]}\textsuperscript{]}\endgroup\ steht, dass
\[
\xi^q-(x^q-x^{aq}+x^{aaq}-x^{a^3q}+x^{a^4q}-\text{etc.}-x^{a^{p-2}q})
\]
\noindent durch $q$ theilbar sein werde, oder von der Form $q\mathbf X$, so dass $\mathbf X$ eine Function mit lauter ganzen Coefficienten bedeutet. Nun ist $q$ nach dem Modulus $p$ einer Potenz von $a$ congruent, es sei $q\equiv a^\mu$. Es ist also
\begingroup\small
\[
\begin{aligned}
&x^q-x^{aq}+x^{aaq}-\text{etc.}-x^{a^{p-2}q}\\
&\qquad-(x^{a^u}-x^{a^{u+1}}+x^{a^{\mu+2}}-x^{a^{\mu+3}}+\text{etc.}-x^{a^{p+\mu-2}})
\end{aligned}
\]
\endgroup
\noindent durch $1-x^p$ theilbar; durch dieselbe Grösse wird aber auch
\[
x^{a^u}-x^{a^{\mu+1}}+x^{a^{\mu+2}}-x^{a^{\mu+3}}+\text{etc.}-x^{a^{p+\mu-2}}\mp\xi
\]
\noindent theilbar seyn, wo das obere oder untere Zeichen gilt, je nachdem $\mu$ gerade oder ungerade, d. i. je nachdem $q$ ein quadratischer Rest oder Nichtrest von $p$ ist. Um durch die doppelten Zeichen keine Verwirrung zu veranlassen, wollen wir annehmen, dass der Buchstab $\varepsilon$ im erstern Falle die Grösse $+1$, im zweiten die Grösse $-1$ bedeutet; es wird folglich
\[
x^q-x^{aq}+x^{aaq}-x^{a^3q}+\text{etc.}-x^{a^{p-2}q}-\varepsilon\xi
\]
\clearpage

\SourceLeaf{038}\ArithmeticHead{31}{R}{FÜNFTER BEWEIS DES FUNDAMENTALSATZES.}
\noindent oder
\[
\xi^q-q\mathbf X-\varepsilon\xi
\]
\noindent durch $1-x^p$, folglich da nach III) $\xi^q\mp\xi p^{\frac12(q-1)}$ durch $\frac{1-x^p}{1-x}$ theilbar ist, auch
\[
\mp\xi p^{\frac12(q-1)}+q\mathbf X+\varepsilon\xi
\]
\noindent oder
\[
q\mathbf X+(\varepsilon\mp p^{\frac12(q-1)})\xi
\]
\noindent durch $\frac{1-x^p}{1-x}$ theilbar seyn. Durch dieselbe Function ist also auch
\[
q\xi\mathbf X+(\varepsilon-(\pm p)^{\frac12q-1})\xi\xi
\]
\noindent und folglich auch
\[
q\xi\mathbf X\pm p[\varepsilon-(\pm p)^{\frac12q-1}]
\]
\noindent theilbar, wo die obern Zeichen sich auf die Form $4m+1$, die untern auf die Form $4m-1$ von $p$ beziehen. Nehmen wir nun an, dass $\xi\mathbf X$ mit $\frac{1-x^p}{1-x}$ auf gewöhnliche Art dividirt den Rest
\[
ax^{p-2}+bx^{p-3}+cx^{p-4}+\text{etc.}+l=R
\]
\noindent gebe, so wird offenbar auch
\[
qR\pm p(\varepsilon-(\pm p)^{\frac12q-1})
\]
\noindent durch $\frac{x^p-1}{x-1}$ theilbar seyn, welches nicht anders möglich ist, als wenn jene Grösse identisch $=0$ wird. Wir haben also
\[
lq\pm p\,(\varepsilon-(\pm p)^{\frac12(q-1)})=0,
\]
\noindent folglich ist
\[
p(\varepsilon-(\pm p)^{\frac12q-1})
\]
\noindent durch $q$ theilbar, mithin auch $\varepsilon-(\pm p)^{\frac12q-1}$, welche Grösse auch durch $\varepsilon\mp p^{\frac12(q-1)}$ dargestellt werden kann, wenn das obere Zeichen genommen wird in den Falle, da wenigstens eine der Grössen $p,q$ von der Form $4m+1$ ist, das untere hingegen, wenn beide die Form $4m+3$ haben. Hieraus leiten wir also die Schlussfolge ab:
\clearpage

\SourceLeaf{039}\ArithmeticHead{32}{L}{ARITHMETIK.\quad NACHLASS.}
\par\indent Ist $q$ ein quadratischer Rest von $p$, so wird, nach der eben ausgesprochenen Bestimmung der Zeichen,
\[
p^{\frac12(q-1)}\equiv\pm1\ (\mathrm{mod.}\ q),
\]
\noindent also auch $p$ quadratischer Rest von $q$, dafern wenigstens eine der Primzahlen $p,q$ von der Form $4m+1$ ist, und Nichtrest von $q$, wenn beide die Form $4m+3$ haben.

Ist hingegen $q$ ein quadratischer Nichtrest von $p$, also $\varepsilon=-1$, so wird
\[
p^{\frac12(q-1)}\equiv\mp1\ (\mathrm{mod.}\ q),
\]
\noindent folglich auch $p$ Nichtrest von $q$ im erstern Falle und Rest im andern.

Dieses ist eben das Fundamentaltheorem selbst.
\ClosingRule
\bigskip
\begin{center}\textbf{BEMERKUNG.}\end{center}
\begingroup\small
\par\indent Die vorstehende Abhandlung, die aus dem September des Jahres 1813 stammt, ist eine von dem lateinischen Texte der \textit{Demonstratio sexta} (1818, Werke II, S. 55 ff.) nur wenig abweichende deutsche Darstellung dieses Beweises des Reziprozitätssatzes; die Zählung der Beweise ist hier umgekehrt, wie in der veröffentlichten Abhandlung. Vergl. den Artikel 20 des \textsc{Bachmann}schen Aufsatzes »Über \textsc{Gauss}' zahlentheoretische Arbeiten«.
\begin{flushright}\textsc{Schlesinger.}\end{flushright}
\ClosingRule\clearpage
\endgroup
